\section{Revisão Bibliográfica}

O artigo publicado por T. Kichdoerfer e M. Ortiz em 2016, \cite{Kirchdoerfer2016}, foi pioneiro nesta temática, sendo intitulada, pelos próprios autores, de \textit{data-driven computing}. No contexto do referido trabalho, os autores analisam a efetividade da abordagem sob a ótica de dois distintos problemas advindos da mecânica dos sólidos. No primeiro, tratam de um problema de equilíbrio estático de treliças tridimensionais não lineares; no segundo, de um problema de elasticidade linear. Para ambos os problemas analisados, os autores observaram boas propriedades de convergência com respeito ao número de pontos de dados e, não obstante, à atribuição de dados locais. Ademais, demonstrou-se que, a medida em que o banco de dados aproxima-se cada vez mais da lei constitutiva (ou material) associada ao problema em análise, a solução advinda do modelo \textit{data-driven} converge para a solução clássica do mesmo. Por fim, os autores ressaltam a robustez da abordagem no que diz respeito as discretizações espaciais empregadas na malha de elementos finitos; em particular, mostrou-se que as soluções de elementos finitos, baseadas no modelo \textit{data-driven}, para o problema de elasticidade linear convergem conjuntamente com respeito ao tamanho da malha e a aproximação pelo banco de dados.   